\documentclass{beamer}
\usetheme{default}
\beamertemplatenavigationsymbolsempty
\usepackage{rembeamer}

\begin{document}

\begin{frame}
  \begin{itemize}
    \item[] 
      Use a histogram to visualize a variable and describe the shape
      of the data (skewness, modality, outliers).  \vspace{0.5cm}
  \end{itemize}
\end{frame}

\begin{frame}
  \frametitle{Histogram}
  \centering
  \begin{tikzpicture}
    \begin{axis}[
      xmin=5,
      xmax=25,
      xlabel={Interest Rate (Percentage)},
      ylabel={Count}
    ]
      \addplot+ [
        hist={bins=8},
        fill=orange,
        draw=black,
        mark=none
      ]
      table [y=interest_rate,col sep=tab] {../data/loans50.dat};
    \end{axis}
  \end{tikzpicture}
\end{frame}

\begin{frame}
  \frametitle{Left Tail}
  \centering
  \begin{tikzpicture}
    \begin{axis}[
      xmin=0,
      xmax=14,
      ylabel={Count}
    ]
      \addplot+ [
        hist={bins=7},
        fill=orange,
        draw=black,
        mark=none
      ]
      table [y=left,col sep=space] {../data/empty.dat};
    \end{axis}
  \end{tikzpicture}
\end{frame}

\begin{frame}
  \frametitle{Right Tail}
  \centering
  \begin{tikzpicture}
    \begin{axis}[
      xmin=0,
      xmax=14,
      ylabel={Count}
    ]
      \addplot+ [
        hist={bins=7},
        fill=orange,
        draw=black,
        mark=none
      ]
      table [y=right,col sep=space] {../data/empty.dat};
    \end{axis}
  \end{tikzpicture}
\end{frame}

\begin{frame}
  \frametitle{Centre - Unimodal}
  \centering
  \begin{tikzpicture}
    \begin{axis}[
      xmin=0,
      xmax=14,
      ylabel={Count}
    ]
      \addplot+ [
        hist={bins=7},
        fill=orange,
        draw=black,
        mark=none
      ]
      table [y=centre,col sep=space] {../data/empty.dat};
    \end{axis}
  \end{tikzpicture}
\end{frame}

\begin{frame}
  \frametitle{Bimodal}
  \centering
  \begin{tikzpicture}
    \begin{axis}[
      xmin=0,
      xmax=14,
      ylabel={Count}
    ]
      \addplot+ [
        hist={bins=7},
        fill=orange,
        draw=black,
        mark=none
      ]
      table [y=bimodal,col sep=space] {../data/empty.dat};
    \end{axis}
  \end{tikzpicture}
\end{frame}

\begin{frame}
  \frametitle{Multimodal}
  \centering
  \begin{tikzpicture}
    \begin{axis}[
      xmin=0,
      xmax=14,
      ylabel={Count}
    ]
      \addplot+ [
        hist={bins=7},
        fill=orange,
        draw=black,
        mark=none
      ]
      table [y=multimodal,col sep=space] {../data/empty.dat};
    \end{axis}
  \end{tikzpicture}
\end{frame}

\begin{frame}
  \frametitle{Outlier}
  \centering
  \begin{tikzpicture}
    \begin{axis}[
      xmin=0,
      xmax=14,
      ylabel={Count}
    ]
      \addplot+ [
        hist={bins=7},
        fill=orange,
        draw=black,
        mark=none
      ]
      table [y=outlier,col sep=space] {../data/empty.dat};
    \end{axis}
  \end{tikzpicture}
\end{frame}

\begin{frame}
  \frametitle{Population Outlier}
  \centering
  \begin{tikzpicture}
    \begin{axis}[
      xmin=0,
      xmax=6,
      ylabel={Count},
      xlabel={Population in Millions}
      ]
      \addplot+ [
        hist={bins=10},
        fill=orange,
        draw=black,
        mark=none
      ]
      table [y=data,col sep=space] {../data/skew_population.dat};
    \end{axis}
  \end{tikzpicture}
\end{frame}

\begin{frame}
  \frametitle{Population Outlier - Logarithm}
  \centering
  \begin{tikzpicture}
    \begin{axis}[
      xmin=-1,
      xmax=1,
      ylabel={Count},
      xlabel={$\log_{10}$(Population in Millions)}
      ]
      \addplot+ [
        hist={bins=10},
        fill=orange,
        draw=black,
        mark=none
      ]
      table [y expr=log10(\thisrow{data}),col sep=space] 
      {../data/skew_population.dat};
    \end{axis}
  \end{tikzpicture}
\end{frame}

\end{document}
